% Chinese questions and concise solutions. Compile using XeLaTeX.
\documentclass[UTF8,10pt,letterpaper,fontset=fandol]{ctexart}
\usepackage[margin=1in]{geometry}
\usepackage{amsmath,amssymb}
\usepackage{enumitem}
\usepackage{needspace}
\usepackage{etoolbox}
\usepackage[unicode,hidelinks]{hyperref}
\pagestyle{plain}
\raggedbottom
\setlength{\parindent}{0pt}
\setlength{\parskip}{1pt}
\setlength{\emergencystretch}{2em}
\setlist[enumerate]{label=(\arabic*),leftmargin=2em,itemsep=0.3em,topsep=0.2em}
\setcounter{secnumdepth}{0}
\ctexset{section={format=\Large\bfseries,beforeskip=1.2ex,afterskip=0.8ex}}
\newcommand{\R}{\mathbb R}
\newcommand{\C}{\mathbb C}
\newcommand{\F}{\mathbb F}
\newcommand{\Span}{\operatorname{span}}
\newcommand{\Solution}{\par\Needspace{5\baselineskip}\textbf{解答}\par}
\newcommand{\CompactMath}{%
  \setlength{\abovedisplayskip}{5pt plus 2pt minus 1pt}%
  \setlength{\belowdisplayskip}{5pt plus 2pt minus 1pt}%
  \setlength{\abovedisplayshortskip}{3pt plus 1pt}%
  \setlength{\belowdisplayshortskip}{5pt plus 2pt minus 1pt}}
\apptocmd{\normalsize}{\CompactMath}{}{}
\begin{document}
\begin{center}
{\LARGE 作业 2}\par
\vspace{0.6em}
2026 年 10 月 10 日\par
姓名：YuanHao Tang\quad 学号：3120265621
\end{center}
\begin{center}\emph{部分题目选自教材中的例题。}\end{center}

\section{第 1 题}
设 \(\mathcal A\in\mathcal L(\mathcal C_2(t)\to\mathcal C_2(t))\)。关于基
\(\alpha_1(t),\alpha_2(t),\alpha_3(t)\)，其节点值为
\[
\alpha_1=\begin{bmatrix}1\\1\\1\end{bmatrix},\quad
\alpha_2=\begin{bmatrix}1\\1\\0\end{bmatrix},\quad
\alpha_3=\begin{bmatrix}1\\0\\0\end{bmatrix},
\]
\(\mathcal A\) 的矩阵表示为
\[
A=\begin{bmatrix}1&1&-1\\0&1&2\\1&2&1\end{bmatrix}.
\]
求 \(\mathcal N(\mathcal A)\) 与 \(\mathcal R(\mathcal A)\)。

\Solution
由 \(Ax=0\)，得 \(x_2=-2x_3\)、\(x_1=3x_3\)。矩阵前两列线性无关，且第三列等于
\(-3\) 倍第一列加 \(2\) 倍第二列。因此
\[
\boxed{\mathcal N(\mathcal A)=\Span\{3\alpha_1(t)-2\alpha_2(t)+\alpha_3(t)\}},
\]
\[
\boxed{\mathcal R(\mathcal A)=\Span\{\alpha_1(t)+\alpha_3(t),\,
\alpha_1(t)+\alpha_2(t)+2\alpha_3(t)\}}.
\]
核与像的维数分别为 \(1\) 与 \(2\)。

\section{第 2 题}
证明下列命题：
\begin{enumerate}
\item 若 \(A\in\F^{n\times n}\)，则
\(S=\{x(t):\dot x(t)=Ax(t)\}\) 是线性空间，且与 \(\C^n\) 同构，其中
\[
\dot x(t)\triangleq
\begin{bmatrix}\dfrac{dx_1(t)}{dt}&\dfrac{dx_2(t)}{dt}&\cdots&\dfrac{dx_n(t)}{dt}\end{bmatrix}^{T}.
\]
\item 给定 \(A\in\C^{n\times n}\)、\(B\in\C^{m\times m}\)，定义
\(\mathcal T_{A,B}:\C^{n\times m}\to\C^{n\times m}\) 为
\[
\mathcal T_{A,B}(X)=AX+XB,
\]
则 \(\mathcal T\in\mathcal L(\C^{n\times m}\to\C^{n\times m})\)。
\end{enumerate}

\Solution
\begin{enumerate}
\item 取 \(\F=\C\)，与题中的 \(\C^n\) 一致；若取 \(\F=\R\)，相应结论为与 \(\R^n\) 同构。
对定义在同一含 \(0\) 区间上的解 \(x,y\) 及标量 \(a,b\)，
\[
\frac{d}{dt}(ax+by)=aAx+bAy=A(ax+by).
\]
零函数也在 \(S\) 中，故 \(S\) 是线性空间。常系数线性微分方程的初值解唯一存在，
所以线性映射 \(E:S\to\C^n\)、\(E(x)=x(0)\) 是双射，其逆为
\(c\mapsto x(t)=e^{At}c\)。因此 \(\boxed{S\cong\C^n}\)。
\item 对任意 \(X,Y\in\C^{n\times m}\)、\(a,b\in\C\)，
\[
\begin{aligned}
\mathcal T_{A,B}(aX+bY)
&=A(aX+bY)+(aX+bY)B\\
&=a(AX+XB)+b(AY+YB)\\
&=a\mathcal T_{A,B}(X)+b\mathcal T_{A,B}(Y).
\end{aligned}
\]
故 \(\mathcal T_{A,B}\) 为线性映射。
\end{enumerate}

\section{第 3 题}
给定 \(\mathcal V^3\) 的两组基：
\[
\begin{array}{ll}
\epsilon_1=(1,0,1)^T,&\eta_1=(1,2,-1)^T,\\
\epsilon_2=(2,1,0)^T,&\eta_2=(2,2,-1)^T,\\
\epsilon_3=(1,1,1)^T,&\eta_3=(2,-1,-1)^T.
\end{array}
\]
定义线性变换 \(\mathcal A\epsilon_i=\eta_i\)，\(i=1,2,3\)。
\begin{enumerate}
\item 求从基 \(\epsilon_1,\epsilon_2,\epsilon_3\) 到基 \(\eta_1,\eta_2,\eta_3\) 的过渡矩阵；
\item 求 \(\mathcal A\) 关于基 \(\epsilon_1,\epsilon_2,\epsilon_3\) 的矩阵表示；
\item 求 \(\mathcal A\) 关于基 \(\eta_1,\eta_2,\eta_3\) 的矩阵表示。
\end{enumerate}

\Solution
\begin{enumerate}
\item 记 \(E=[\epsilon_1\ \epsilon_2\ \epsilon_3]\)、\(H=[\eta_1\ \eta_2\ \eta_3]\)，
采用 \(H=ET\) 的过渡矩阵约定，得
\[
\boxed{T=E^{-1}H=\frac12\begin{bmatrix}-4&-3&3\\2&3&3\\2&1&-5\end{bmatrix}}.
\]
\item 由 \(\mathcal A(E)=H=ET\)，得 \(\boxed{[\mathcal A]_{\epsilon}=T}\)。
\item 换基后 \(\boxed{[\mathcal A]_{\eta}=T^{-1}[\mathcal A]_{\epsilon}T=T}\)。
\end{enumerate}

\section{第 4 题}
记 \(\R_n[\lambda]\) 为所有实系数、次数不超过 \(n\) 的多项式构成的集合。
定义线性变换 \(\mathcal A\in\mathcal L(\R_n[\lambda]\mapsto\R_n[\lambda])\) 为
\[
\mathcal A(f(\lambda))=\frac{d}{d\lambda}f(\lambda),\qquad f(\lambda)\in\R_n[\lambda],
\]
其中 \(\mathcal L(\cdot)\) 表示线性变换的集合。当 \(n=2\) 时，求：
\begin{enumerate}
\item \(\mathcal A\) 关于基 \(\alpha_1(\lambda)=1+\lambda+\lambda^2\)、
\(\alpha_2(\lambda)=1+\lambda\)、\(\alpha_3(\lambda)=1+2\lambda\) 的矩阵表示；
\item \(\mathcal A\) 的核 \(\mathcal N(\mathcal A)\) 与像 \(\mathcal R(\mathcal A)\)。
\end{enumerate}

\Solution
\begin{enumerate}
\item 由 \(1=2\alpha_2-\alpha_3\)，得
\[
\mathcal A\alpha_1=\alpha_3,\quad
\mathcal A\alpha_2=2\alpha_2-\alpha_3,\quad
\mathcal A\alpha_3=4\alpha_2-2\alpha_3.
\]
将各像的坐标按列排列，得
\[
\boxed{[\mathcal A]_{\alpha}=\begin{bmatrix}0&0&0\\0&2&4\\1&-1&-2\end{bmatrix}}.
\]
\item 导数为零的多项式为常数，多项式的导数遍及 \(\R_1[\lambda]\)。故
\[
\boxed{\mathcal N(\mathcal A)=\Span\{1\}},\qquad
\boxed{\mathcal R(\mathcal A)=\Span\{1,\lambda\}=\R_1[\lambda]}.
\]
核与像的维数分别为 \(1\) 与 \(2\)。
\end{enumerate}

\section{第 5 题}
设 \(\mathcal A\) 是四维线性空间 \(\mathcal V\) 到自身的线性变换，
\(\alpha_1,\alpha_2,\alpha_3,\alpha_4\) 为 \(\mathcal V\) 的一组基。
\(\mathcal A\) 关于这组基的矩阵表示为
\[
A=\begin{bmatrix}
5&-2&-4&3\\3&-1&-3&2\\-3&\tfrac12&\tfrac92&-\tfrac52\\-10&3&11&-7
\end{bmatrix}.
\]
\begin{enumerate}
\item 求 \(\mathcal A\) 关于下列基的矩阵表示：
\[
\begin{aligned}
\eta_1&=\alpha_1+2\alpha_2+\alpha_3+\alpha_4,\\
\eta_2&=2\alpha_1+3\alpha_2+\alpha_3,\\
\eta_3&=\alpha_3,\qquad \eta_4=\alpha_4.
\end{aligned}
\]
\item 求 \(\mathcal A\) 的核 \(\mathcal N(\mathcal A)\) 与像 \(\mathcal R(\mathcal A)\)。
\end{enumerate}

\Solution
\begin{enumerate}
\item 由 \([\eta_1\ \eta_2\ \eta_3\ \eta_4]=[\alpha_1\ \alpha_2\ \alpha_3\ \alpha_4]P\)，
\[
P=\begin{bmatrix}1&2&0&0\\2&3&0&0\\1&1&1&0\\1&0&0&1\end{bmatrix},\qquad
\boxed{[\mathcal A]_{\eta}=P^{-1}AP=
\begin{bmatrix}0&0&6&-5\\0&0&-5&4\\0&0&\tfrac72&-\tfrac32\\0&0&5&-2\end{bmatrix}}.
\]
\item 新矩阵前两列为零，后两列线性无关，故核与像的维数均为 \(2\)，且
\[
\boxed{\mathcal N(\mathcal A)=\Span\{\eta_1,\eta_2\}},\qquad
\boxed{\mathcal R(\mathcal A)=\Span\{r_1,r_2\}}.
\]
其中，由 \(\eta_3=\alpha_3\)、\(\eta_4=\alpha_4\)，可取
\[
\begin{aligned}
r_1&=\mathcal A\alpha_3=-4\alpha_1-3\alpha_2+\tfrac92\alpha_3+11\alpha_4,\\
r_2&=\mathcal A\alpha_4=3\alpha_1+2\alpha_2-\tfrac52\alpha_3-7\alpha_4.
\end{aligned}
\]
\end{enumerate}
\end{document}
